\section{Паросочетание: Хопкрофт--Карп}
\texttt{graph[u]} содержит номера соседей левой вершины среди \texttt{rightCount} правых вершин; возвращает массивы партнёров максимального паросочетания слева и справа, $-1$ означает свободную вершину.
Доли нумеруются с нуля независимо; пустые доли и кратные рёбра допустимы, $n+1<\texttt{INT\_MAX}$.
Время $O((V+E)\sqrt{V+1})$, дополнительная память $O(V)$, нужен стек глубины $O(V)$.
\begin{lstlisting}[name=advanced_graph_hopcroft_karp]
pair<vector<int>, vector<int>> hopcroft_karp(
        const vector<vector<int>>& graph, int rightCount) {
    int n = graph.size();
    vector<int> left(n, -1), right(rightCount, -1), distance(n), position(n);
    while (true) {
        queue<int> queue;
        fill(distance.begin(), distance.end(), -1);
        fill(position.begin(), position.end(), 0);
        for (int u = 0; u < n; ++u)
            if (left[u] == -1) { distance[u] = 0; queue.push(u); }
        int shortest = n + 1;
        while (!queue.empty()) {
            int u = queue.front(); queue.pop();
            if (distance[u] >= shortest) continue;
            for (int v : graph[u]) {
                int next = right[v];
                if (next == -1) shortest = distance[u] + 1;
                else if (distance[next] == -1) {
                    distance[next] = distance[u] + 1;
                    queue.push(next);
                }
            }
        }
        if (shortest == n + 1) break;
        auto augment = [&](auto&& self, int u) -> bool {
            for (int& i = position[u]; i < (int)graph[u].size(); ++i) {
                int v = graph[u][i], next = right[v];
                if ((next == -1 && distance[u] + 1 == shortest) ||
                    (next != -1 && distance[next] == distance[u] + 1 &&
                     self(self, next))) {
                    left[u] = v; right[v] = u;
                    return true;
                }
            }
            distance[u] = -1;
            return false;
        };
        for (int u = 0; u < n; ++u)
            if (left[u] == -1) augment(augment, u);
    }
    return {left, right};
}
\end{lstlisting}
% Максимальное паросочетание двудольного графа: \texttt{graph[u]} содержит
% соседей левой вершины среди \texttt{rightCount} правых вершин.
% Результат --- пара массивов партнёров слева и справа, $-1$ означает свободную
% вершину. Пустые доли и кратные рёбра допустимы. Время
% $O((V+E)\sqrt{V+1})$, память $O(V)$ сверх входа.
% Все восемь реализаций ниже авторские, не дословные заимствования.
% Источники алгоритмов: J. Hopcroft, R. Karp, ``An $n^{5/2}$ Algorithm for
% Maximum Matchings in Bipartite Graphs'' (1973); H. Kuhn, ``The Hungarian
% Method for the Assignment Problem'' (1955); стандартные алгоритмы
% последовательных кратчайших путей, Диница, Тарьяна и Краскала.
% Общие предпосылки: вне листингов подключён \texttt{bits/stdc++.h}, заданы
% \texttt{using namespace std; using ll = long long; using db = double;}.
% Индексы, размеры, степени и служебные размеры помещаются в \texttt{int};
% вершины нумеруются с нуля, концы рёбер корректны. Веса, ёмкости и суммы
% имеют тип \texttt{ll}. Все выполняемые арифметические операции обязаны
% помещаться в свой тип; дополнительные ограничения приведены ниже.
% Рекурсивные лямбды требуют стека глубины $O(V)$ (в динамической связности
% $O(\log Q)$). Ошибки входа, кроме явно проверяемых, нарушают контракт.
%
% По чередующимся путям из свободных левых вершин восстанавливается
% минимальное вершинное покрытие (теорема Кёнига). Итеративный DFS снимает
% ограничение на глубину стека. Здесь дополнительно требуется $n+1<\texttt{INT\_MAX}$.

\section{Прямоугольное назначение: венгерский алгоритм}
Для прямоугольной матрицы стоимостей $n\times m$ возвращает минимальную стоимость назначения каждой строки разному столбцу и номера выбранных столбцов.
При $n>m$ возвращает \texttt{nullopt}, пустой вход даёт нулевую стоимость; все назначения разрешены, отрицательные стоимости допустимы.
Требуется $n+1,m+1<\texttt{INT\_MAX}$; вся арифметика помещается в \texttt{ll}, конечные \texttt{slack} и \texttt{delta} меньше \texttt{LLONG\_MAX}.
Время $O(n^2m)$, дополнительная память $O(n+m)$.
\begin{lstlisting}[name=advanced_graph_hungarian]
optional<pair<ll, vector<int>>> hungarian(const vector<vector<ll>>& cost) {
    int n = cost.size(), m = n == 0 ? 0 : cost[0].size();
    if (n > m) return nullopt;
    vector<ll> rowPotential(n + 1), columnPotential(m + 1);
    vector<int> matched(m + 1), previous(m + 1);
    for (int row = 1; row <= n; ++row) {
        matched[0] = row;
        int column = 0;
        vector<ll> slack(m + 1, LLONG_MAX);
        vector<char> used(m + 1);
        do {
            used[column] = true;
            int currentRow = matched[column], nextColumn = -1;
            ll delta = LLONG_MAX;
            for (int j = 1; j <= m; ++j) if (!used[j]) {
                ll reduced = cost[currentRow - 1][j - 1]
                           - rowPotential[currentRow] - columnPotential[j];
                if (reduced < slack[j]) {
                    slack[j] = reduced;
                    previous[j] = column;
                }
                if (slack[j] < delta) { delta = slack[j]; nextColumn = j; }
            }
            for (int j = 0; j <= m; ++j) {
                if (used[j]) {
                    rowPotential[matched[j]] += delta;
                    columnPotential[j] -= delta;
                } else if (slack[j] != LLONG_MAX) slack[j] -= delta;
            }
            if (nextColumn == -1)
                throw logic_error("no free column");
            column = nextColumn;
        } while (matched[column] != 0);
        do {
            int nextColumn = previous[column];
            matched[column] = matched[nextColumn];
            column = nextColumn;
        } while (column != 0);
    }
    vector<int> assignment(n);
    for (int j = 1; j <= m; ++j)
        if (matched[j] != 0) assignment[matched[j] - 1] = j - 1;
    ll total = 0;
    for (int i = 0; i < n; ++i) total += cost[i][assignment[i]];
    return pair{total, assignment};
}
\end{lstlisting}
% Минимальная стоимость назначения каждой из $n$ строк разному столбцу
% матрицы $n\times m$, $n\le m$. Все назначения разрешены, стоимости могут
% быть отрицательными. Возвращаются стоимость и выбранный столбец каждой
% строки; при $n>m$ --- \texttt{nullopt}, пустая матрица даёт нулевую стоимость.
% Время $O(n^2m)$, память $O(n+m)$ сверх матрицы.
% Матрица прямоугольная. Требуется $n+1,m+1<\texttt{INT\_MAX}$.
% Все приведённые стоимости, потенциалы и промежуточные разности помещаются
% в \texttt{ll}; конечные значения \texttt{slack} и \texttt{delta} строго
% меньше \texttt{LLONG\_MAX}. Итоговая сумма также помещается в \texttt{ll}.
% Бесконечность не заменяет запрещённое назначение.
%
% При $n>m$ можно транспонировать задачу, если требуется покрыть меньшую
% долю. Для разреженных разрешённых назначений удобнее поток минимальной
% стоимости; для максимизации нужно безопасно изменить знак стоимостей.

\section{Поток минимальной стоимости с потенциалами}
Для дуг \texttt{(from,to,capacity,cost)} возвращает поток до \texttt{limit}, его минимальную стоимость и потоки по исходным ID; при недостижимом лимите поток максимальный.
Требуются $s\ne t$, неотрицательные ёмкости и лимит, отсутствие любых отрицательных циклов положительной ёмкости в начальной остаточной сети; достижимый отрицательный цикл вызывает \texttt{domain\_error}.
Петли и кратные дуги допустимы; $2E<\texttt{INT\_MAX}$, стоимость не равна \texttt{LLONG\_MIN}, вся арифметика помещается в \texttt{ll}, конечные расстояния меньше \texttt{LLONG\_MAX}.
Время $O(VE+K(E+V)\log(V+1))$ для $K$ увеличений, память $O(V+E)$.
\begin{lstlisting}[name=advanced_graph_min_cost_flow]
tuple<ll, ll, vector<ll>> min_cost_flow(int n,
        const vector<tuple<int, int, ll, ll>>& input, int s, int t, ll limit) {
    struct Edge { int from, to; ll capacity, cost; };
    vector<Edge> edges;
    vector<vector<int>> graph(n);
    for (auto [u, v, capacity, cost] : input) {
        int id = edges.size();
        graph[u].push_back(id); graph[v].push_back(id + 1);
        edges.push_back({u, v, capacity, cost});
        edges.push_back({v, u, 0, -cost});
    }
    vector<ll> potential(n, LLONG_MAX);
    potential[s] = 0;
    for (int iteration = 0; iteration < n; ++iteration) {
        bool changed = false;
        for (const auto& edge : edges) {
            if (edge.capacity == 0 || potential[edge.from] == LLONG_MAX) continue;
            ll candidate = potential[edge.from] + edge.cost;
            if (candidate < potential[edge.to]) {
                potential[edge.to] = candidate;
                changed = true;
            }
        }
        if (!changed) break;
        if (iteration == n - 1) throw domain_error("reachable negative cost cycle");
    }
    for (ll& value : potential) if (value == LLONG_MAX) value = 0;
    ll flow = 0, total = 0;
    while (flow < limit) {
        vector<ll> distance(n, LLONG_MAX);
        vector<int> previous(n, -1);
        priority_queue<pair<ll, int>, vector<pair<ll, int>>, greater<>> queue;
        distance[s] = 0; queue.push({0, s});
        while (!queue.empty()) {
            auto [current, u] = queue.top(); queue.pop();
            if (current != distance[u]) continue;
            for (int id : graph[u]) {
                const auto& edge = edges[id];
                if (edge.capacity == 0) continue;
                ll reduced = edge.cost + potential[u] - potential[edge.to];
                ll candidate = current + reduced;
                if (candidate < distance[edge.to]) {
                    distance[edge.to] = candidate;
                    previous[edge.to] = id;
                    queue.push({candidate, edge.to});
                }
            }
        }
        if (previous[t] == -1) break;
        for (int v = 0; v < n; ++v)
            if (distance[v] != LLONG_MAX) potential[v] += distance[v];
        ll increase = limit - flow;
        for (int v = t; v != s; v = edges[previous[v]].from)
            increase = min(increase, edges[previous[v]].capacity);
        for (int v = t; v != s; v = edges[previous[v]].from) {
            int id = previous[v];
            total += increase * edges[id].cost;
            edges[id].capacity -= increase;
            edges[id ^ 1].capacity += increase;
        }
        flow += increase;
    }
    vector<ll> edgeFlow(input.size());
    for (int id = 0; id < (int)input.size(); ++id)
        edgeFlow[id] = edges[2 * id + 1].capacity;
    return {flow, total, edgeFlow};
}
\end{lstlisting}
% Ориентированные рёбра заданы кортежами \texttt{(from,to,capacity,cost)}.
% Возвращаются отправленный поток (не более \texttt{limit}), его стоимость
% и потоки по исходным ID. Если лимит недостижим, возвращается максимальный
% возможный поток. Требуются $s\ne t$, неотрицательные ёмкости и лимит.
% Петли и кратные дуги поддержаны. Начальная остаточная сеть не должна
% содержать отрицательных циклов положительной ёмкости.
% В частности, достижимый из $s$ отрицательный цикл явно обнаруживается
% Беллманом--Фордом и вызывает \texttt{domain\_error}; недостижимые циклы
% не проверяются и также исключены контрактом для глобальной оптимальности.
% Время $O(VE+K(E+V)\log(V+1))$ для $K$ увеличений, память $O(V+E)$.
% Требуется $2E<\texttt{INT\_MAX}$; стоимость не равна
% \texttt{LLONG\_MIN}. Все суммы расстояний, потенциалов, приведённые стоимости,
% разности ёмкостей, произведения увеличения на стоимость дуги и накопленные
% стоимости помещаются в \texttt{ll}. Любое конечное расстояние меньше
% \texttt{LLONG\_MAX}; этот маркер никогда не участвует в сложении.
%
% Для произвольных отрицательных циклов нужны отмена циклов или другой
% алгоритм. Проверка из фиктивного источника во все вершины может выявить
% также недостижимые отрицательные циклы. Масштабирование ёмкостей полезно
% для больших значений потока; число увеличений здесь не полиномиально по
% длине двоичной записи ёмкостей.

\section{Допустимая циркуляция с нижними границами}
Для дуг \texttt{(from,to,lower,upper)} возвращает потоки по ID с сохранением баланса каждой вершины либо \texttt{nullopt}; источника и стока нет.
Требуется $0\le lower\le upper$, неверные границы дают \texttt{nullopt}; петли и пустая сеть допустимы.
Требуется $n+2<\texttt{INT\_MAX}$, $2(E+n)<\texttt{INT\_MAX}$; вся арифметика помещается в \texttt{ll}, баланс не равен \texttt{LLONG\_MIN}.
Время $O(V^2(E+V))$, память $O(V+E)$, нужен стек глубины $O(V)$.
\begin{lstlisting}[name=advanced_graph_circulation]
optional<vector<ll>> feasible_circulation(int n,
        const vector<tuple<int, int, ll, ll>>& input) {
    struct Edge { int to; ll capacity; };
    int source = n, sink = n + 1;
    vector<Edge> edges;
    vector<vector<int>> graph(n + 2);
    auto add = [&](int u, int v, ll capacity) {
        int id = edges.size();
        graph[u].push_back(id); graph[v].push_back(id + 1);
        edges.push_back({v, capacity}); edges.push_back({u, 0});
    };
    vector<ll> balance(n);
    for (auto [u, v, lower, upper] : input) {
        if (lower < 0 || lower > upper) return nullopt;
        add(u, v, upper - lower);
        balance[u] -= lower; balance[v] += lower;
    }
    ll required = 0;
    for (int v = 0; v < n; ++v) {
        if (balance[v] > 0) { add(source, v, balance[v]); required += balance[v]; }
        if (balance[v] < 0) add(v, sink, -balance[v]);
    }
    vector<int> level(n + 2), position(n + 2);
    ll flow = 0;
    while (flow < required) {
        fill(level.begin(), level.end(), -1);
        queue<int> queue;
        queue.push(source); level[source] = 0;
        while (!queue.empty()) {
            int u = queue.front(); queue.pop();
            for (int id : graph[u]) {
                int v = edges[id].to;
                if (edges[id].capacity > 0 && level[v] == -1) {
                    level[v] = level[u] + 1; queue.push(v);
                }
            }
        }
        if (level[sink] == -1) return nullopt;
        fill(position.begin(), position.end(), 0);
        auto push = [&](auto&& self, int u, ll amount) -> ll {
            if (u == sink) return amount;
            for (int& i = position[u]; i < (int)graph[u].size(); ++i) {
                int id = graph[u][i], v = edges[id].to;
                if (edges[id].capacity == 0 || level[v] != level[u] + 1) continue;
                ll sent = self(self, v, min(amount, edges[id].capacity));
                if (sent == 0) continue;
                edges[id].capacity -= sent; edges[id ^ 1].capacity += sent;
                return sent;
            }
            return 0;
        };
        while (flow < required) {
            ll sent = push(push, source, required - flow);
            if (sent == 0) break;
            flow += sent;
        }
    }
    vector<ll> result(input.size());
    for (int id = 0; id < (int)input.size(); ++id)
        result[id] = get<2>(input[id]) + edges[2 * id + 1].capacity;
    return result;
}
\end{lstlisting}
% Для каждой дуги \texttt{(from,to,lower,upper)} найти поток в границах,
% сохраняющий баланс в каждой вершине. Возвращается массив потоков по ID
% либо \texttt{nullopt}. Это именно циркуляция без источника и стока.
% Требуется $0\le lower\le upper$; неверные границы дают \texttt{nullopt}.
% Петли допустимы, пустая сеть допустима. Время $O(V^2(E+V))$, память $O(V+E)$.
% Диниц полностью локален. Требуется $n+2<\texttt{INT\_MAX}$,
% $2(E+n)<\texttt{INT\_MAX}$. Балансы, сумма положительных балансов,
% все промежуточные суммы и восстановленные потоки помещаются в \texttt{ll};
% баланс не равен \texttt{LLONG\_MIN}, поскольку может отрицаться.
%
% Для потока из $s$ в $t$ добавляется дуга $t\to s$ с подходящей верхней
% границей; её поток определяет значение исходного потока. Максимизация
% требует отдельной фазы в остаточной сети после удаления служебных дуг.

\section{Точки сочленения и рёберные блоки двусвязности}
Для неориентированного мультиграфа, заданного парами концов рёбер, возвращает флаги точек сочленения, блоки вершинной двусвязности как списки ID рёбер и изолированные вершины.
ID --- позиция ребра во входе; каждый мост и каждая петля образуют отдельный блок, изолированные вершины пустых блоков не создают.
Требуется $2E<\texttt{INT\_MAX}$; время и память $O(V+E)$, нужен стек глубины $O(V)$.
\begin{lstlisting}[name=advanced_graph_biconnected]
tuple<vector<char>, vector<vector<int>>, vector<int>> biconnected_blocks(
        int n, const vector<pair<int, int>>& edges) {
    vector<vector<pair<int, int>>> graph(n);
    vector<vector<int>> blocks;
    for (int id = 0; id < (int)edges.size(); ++id) {
        auto [u, v] = edges[id];
        if (u == v) blocks.push_back({id});
        graph[u].push_back({v, id}); graph[v].push_back({u, id});
    }
    vector<int> tin(n, -1), low(n), stack, isolated;
    vector<char> articulation(n);
    int timer = 0;
    auto dfs = [&](auto&& self, int u, int parentEdge) -> void {
        tin[u] = low[u] = timer++;
        int children = 0;
        for (auto [v, id] : graph[u]) {
            if (id == parentEdge || v == u) continue;
            if (tin[v] == -1) {
                ++children; stack.push_back(id);
                self(self, v, id);
                low[u] = min(low[u], low[v]);
                if (low[v] >= tin[u]) {
                    if (parentEdge != -1) articulation[u] = true;
                    vector<int> block;
                    do { block.push_back(stack.back()); stack.pop_back(); }
                    while (block.back() != id);
                    blocks.push_back(move(block));
                }
            } else if (tin[v] < tin[u]) {
                stack.push_back(id); low[u] = min(low[u], tin[v]);
            }
        }
        if (parentEdge == -1) articulation[u] = children > 1;
    };
    for (int u = 0; u < n; ++u) if (tin[u] == -1) {
        if (graph[u].empty()) isolated.push_back(u);
        dfs(dfs, u, -1);
    }
    return {articulation, blocks, isolated};
}
\end{lstlisting}
% Неориентированный мультиграф задан парами концов; ID равен позиции ребра.
% Возвращаются флаги точек сочленения, блоки как списки ID рёбер и список
% изолированных вершин. Здесь блоки --- компоненты вершинной двусвязности,
% а не компоненты после удаления мостов. Мост образует отдельный блок,
% каждая петля по принятому соглашению тоже отдельный блок. Изолированные
% вершины не создают пустых блоков и перечисляются отдельно.
% Время и память $O(V+E)$.
% Пропускается только родительское ребро по ID, а не все рёбра к родителю:
% это существенно для кратных рёбер. Требуется $2E<\texttt{INT\_MAX}$.
%
% Из блоков можно построить block-cut forest. Для петель соглашение о
% блоках нужно согласовать с потребителем результата; обычные утверждения
% о block-cut tree часто формулируются для графов без петель. Итеративный
% обход позволяет обрабатывать длинные цепочки без большого системного стека.

\section{Офлайн-динамическая связность}
Операция \texttt{(type,u,v)}: 0 добавляет неориентированное ребро, 1 удаляет одну копию, 2 спрашивает связность; возвращает ответы $0/1$ в порядке запросов.
Петли допустимы, копии удаляются в порядке LIFO; удаление отсутствующего ребра и неизвестный тип вызывают \texttt{invalid\_argument}.
Требуется $4\max(1,Q)<\texttt{INT\_MAX}$, число записей интервалов и размер истории помещаются в \texttt{int}.
Время $O(Q\log(Q+1)\log(V+1))$, память $O(V+Q\log(Q+1))$, глубина рекурсии $O(\log(Q+1))$.
\begin{lstlisting}[name=advanced_graph_dynamic_connectivity]
vector<int> offline_dynamic_connectivity(int n,
        const vector<tuple<int, int, int>>& operations) {
    int q = operations.size();
    if (q == 0) return {};
    vector<vector<pair<int, int>>> tree(4 * q);
    auto add = [&](auto&& self, int node, int l, int r, int a, int b,
                   pair<int, int> edge) -> void {
        if (a >= r || b <= l) return;
        if (a <= l && r <= b) { tree[node].push_back(edge); return; }
        int middle = l + (r - l) / 2;
        self(self, node * 2, l, middle, a, b, edge);
        self(self, node * 2 + 1, middle, r, a, b, edge);
    };
    map<pair<int, int>, vector<int>> active;
    for (int time = 0; time < q; ++time) {
        auto [type, u, v] = operations[time];
        if (u > v) swap(u, v);
        pair<int, int> edge = {u, v};
        if (type == 0) active[edge].push_back(time);
        else if (type == 1) {
            auto it = active.find(edge);
            if (it == active.end() || it->second.empty())
                throw invalid_argument("removing inactive edge");
            int start = it->second.back(); it->second.pop_back();
            add(add, 1, 0, q, start, time, edge);
        } else if (type != 2) throw invalid_argument("unknown operation type");
    }
    for (const auto& [edge, starts] : active)
        for (int start : starts) add(add, 1, 0, q, start, q, edge);
    vector<int> parent(n), size(n, 1), answers;
    iota(parent.begin(), parent.end(), 0);
    auto find = [&](int v) {
        while (parent[v] != v) v = parent[v];
        return v;
    };
    vector<pair<int, int>> history;
    auto solve = [&](auto&& self, int node, int l, int r) -> void {
        int snapshot = history.size();
        for (auto [u, v] : tree[node]) {
            u = find(u); v = find(v);
            if (u == v) continue;
            if (size[u] < size[v]) swap(u, v);
            history.push_back({v, size[u]});
            parent[v] = u; size[u] += size[v];
        }
        if (r - l == 1) {
            auto [type, u, v] = operations[l];
            if (type == 2) answers.push_back(find(u) == find(v));
        } else {
            int middle = l + (r - l) / 2;
            self(self, node * 2, l, middle);
            self(self, node * 2 + 1, middle, r);
        }
        while ((int)history.size() > snapshot) {
            auto [v, oldSize] = history.back(); history.pop_back();
            size[parent[v]] = oldSize; parent[v] = v;
        }
    };
    solve(solve, 1, 0, q);
    return answers;
}
\end{lstlisting}
% Операции --- тройки \texttt{(type,u,v)}: $0$ добавляет неориентированное
% ребро, $1$ удаляет одну его копию, $2$ спрашивает связность вершин.
% Возвращаются ответы $0/1$ в порядке запросов. Пары нормализуются; копии
% одинакового ребра удаляются в порядке LIFO, что не меняет связность.
% Удаление отсутствующего ребра и неизвестный тип вызывают
% \texttt{invalid\_argument}. Петли разрешены, вершина связна сама с собой.
% Время $O(Q\log(Q+1)\log(V+1))$, память $O(V+Q\log(Q+1))$.
% DSU без сжатия путей, с объединением по размеру и откатом к снимку.
% Интервал активности полуоткрытый: от добавления до удаления.
% Требуется $4\max(1,Q)<\texttt{INT\_MAX}$, а суммарное число записей
% интервалов и история отката также помещаются в \texttt{int}.
%
% В историю можно добавить число компонент или агрегаты компонент.
% Для онлайн-удалений нужны существенно более сложные структуры; обычная
% DSU со сжатием путей не подходит для этого отката и оценки сложности.

\iffalse
\section{Дерево реконструкции Краскала}
Для неориентированных рёбер \texttt{(u,v,weight)} возвращает лес реконструкции с полями \texttt{children}, \texttt{parent}, \texttt{weight}, \texttt{mergeEdge}, \texttt{roots}; исходные вершины --- листья $[0,n)$.
Узел объединения хранит двух детей, вес и исходный ID ребра; у листьев \texttt{mergeEdge} равен $-1$, вес 0 условен, у корней родитель $-1$.
Вес LCA различных листьев одной компоненты --- минимально возможный максимум веса на пути между ними; петли игнорируются, отрицательные веса допустимы, пустой граф даёт пустой лес.
Требуется $2n<\texttt{INT\_MAX}$; время $O(E\log(E+1)+V+E\alpha(V))$, память $O(V+E)$.
\begin{lstlisting}[name=advanced_graph_kruskal_tree]
auto kruskal_reconstruction_tree(int n,
        const vector<tuple<int, int, ll>>& edges) {
    struct Result {
        vector<array<int, 2>> children;
        vector<int> parent;
        vector<ll> weight;
        vector<int> mergeEdge, roots;
    };
    Result result{vector<array<int, 2>>(n, {-1, -1}), vector<int>(n, -1),
                  vector<ll>(n), vector<int>(n, -1), {}};
    vector<int> dsu(n), size(n, 1), top(n), order(edges.size());
    iota(dsu.begin(), dsu.end(), 0);
    iota(top.begin(), top.end(), 0);
    iota(order.begin(), order.end(), 0);
    auto find = [&](int v) {
        while (v != dsu[v]) { dsu[v] = dsu[dsu[v]]; v = dsu[v]; }
        return v;
    };
    sort(order.begin(), order.end(), [&](int a, int b) {
        return pair{get<2>(edges[a]), a} < pair{get<2>(edges[b]), b};
    });
    for (int id : order) {
        auto [u, v, weight] = edges[id];
        u = find(u); v = find(v);
        if (u == v) continue;
        int node = result.parent.size();
        result.children.push_back({top[u], top[v]});
        result.parent[top[u]] = result.parent[top[v]] = node;
        result.parent.push_back(-1); result.weight.push_back(weight);
        result.mergeEdge.push_back(id);
        if (size[u] < size[v]) swap(u, v);
        dsu[v] = u; size[u] += size[v]; top[u] = node;
    }
    for (int v = 0; v < n; ++v)
        if (dsu[v] == v) result.roots.push_back(top[v]);
    return result;
}
\end{lstlisting}
% Неориентированные рёбра --- \texttt{(u,v,weight)}. Возвращается локальный
% тип с полями \texttt{children}, \texttt{parent}, \texttt{weight},
% \texttt{mergeEdge}, \texttt{roots}. Исходные вершины --- листья $[0,n)$;
% каждое успешное объединение создаёт новую вершину с двумя детьми и весом
% объединяющего ребра. \texttt{mergeEdge} хранит его исходный ID, для листа
% равен $-1$; вес листа $0$ условен и не участвует в сравнениях.
% Для несвязного графа получается лес, корни имеют родителя $-1$.
% Вес LCA двух различных листьев одной компоненты --- минимально возможный
% максимум веса ребра на соединяющем пути. Петли игнорируются, равные веса
% обрабатываются по ID. Время $O(E\log(E+1)+V+E\alpha(V))$, память $O(V+E)$.
% Требуется $2n<\texttt{INT\_MAX}$. Отрицательные веса допустимы;
% арифметика над весами не выполняется. Пустой граф даёт пустой лес.
%
% Двоичные подъёмы или Euler tour с RMQ дают быстрые LCA-запросы.
% Для запросов между разными компонентами нужно отдельно проверять корень;
% для запроса вершины к самой себе порог следует определить соглашением.

\fi
\iffalse
\section{Второй остов: другой набор рёбер}
Для неориентированного взвешенного мультиграфа возвращает \texttt{mstWeight}, ID рёбер \texttt{mstEdges} и \texttt{alternative}: стоимость другого остова, ID добавляемого и удаляемого ребра.
При несвязности или $n=0$ результат отсутствует; при единственном остове отсутствует альтернатива, равная стоимость разрешена --- отличается набор ID.
Все промежуточные суммы и разности весов помещаются в \texttt{ll}, размеры таблиц --- в \texttt{int}; отрицательные веса, включая \texttt{LLONG\_MIN}, допустимы.
Время $O(E\log(E+1)+(V+E)\log(V+1))$, память $O(E+V\log(V+1))$, нужен стек глубины $O(V)$.
\begin{lstlisting}[name=advanced_graph_second_mst]
auto second_spanning_tree(int n, const vector<tuple<int, int, ll>>& edges) {
    struct Result {
        ll mstWeight = 0;
        vector<int> mstEdges;
        optional<tuple<ll, int, int>> alternative;
    };
    if (n == 0) return optional<Result>{};
    int m = edges.size();
    vector<int> parent(n), size(n, 1), order(m);
    iota(parent.begin(), parent.end(), 0);
    iota(order.begin(), order.end(), 0);
    auto find = [&](int v) {
        while (v != parent[v]) { parent[v] = parent[parent[v]]; v = parent[v]; }
        return v;
    };
    sort(order.begin(), order.end(), [&](int a, int b) {
        return pair{get<2>(edges[a]), a} < pair{get<2>(edges[b]), b};
    });
    Result result;
    vector<char> chosen(m);
    vector<vector<pair<int, int>>> graph(n);
    for (int id : order) {
        auto [u, v, weight] = edges[id];
        int a = find(u), b = find(v);
        if (a == b) continue;
        if (size[a] < size[b]) swap(a, b);
        parent[b] = a; size[a] += size[b];
        result.mstWeight += weight; result.mstEdges.push_back(id); chosen[id] = true;
        graph[u].push_back({v, id}); graph[v].push_back({u, id});
    }
    if ((int)result.mstEdges.size() != n - 1) return optional<Result>{};
    int levels = 1;
    while ((1LL << levels) <= n) ++levels;
    vector<vector<int>> up(levels, vector<int>(n));
    vector<vector<int>> maximum(levels, vector<int>(n, -1));
    vector<int> depth(n, -1), traversal = {0};
    depth[0] = 0;
    for (int i = 0; i < (int)traversal.size(); ++i) {
        int u = traversal[i];
        for (auto [v, id] : graph[u]) if (depth[v] == -1) {
            depth[v] = depth[u] + 1; up[0][v] = u; maximum[0][v] = id;
            traversal.push_back(v);
        }
    }
    auto heavier = [&](int a, int b) {
        if (a == -1) return b;
        if (b == -1) return a;
        return pair{get<2>(edges[a]), a} < pair{get<2>(edges[b]), b} ? b : a;
    };
    for (int k = 1; k < levels; ++k)
        for (int v = 0; v < n; ++v) {
            up[k][v] = up[k - 1][up[k - 1][v]];
            maximum[k][v] = heavier(maximum[k - 1][v],
                                    maximum[k - 1][up[k - 1][v]]);
        }
    auto pathMaximum = [&](int u, int v) {
        int answer = -1;
        if (depth[u] < depth[v]) swap(u, v);
        int difference = depth[u] - depth[v];
        for (int k = 0; k < levels; ++k) if ((difference >> k) & 1) {
            answer = heavier(answer, maximum[k][u]); u = up[k][u];
        }
        if (u == v) return answer;
        for (int k = levels - 1; k >= 0; --k) if (up[k][u] != up[k][v]) {
            answer = heavier(answer, heavier(maximum[k][u], maximum[k][v]));
            u = up[k][u]; v = up[k][v];
        }
        return heavier(answer, heavier(maximum[0][u], maximum[0][v]));
    };
    for (int id = 0; id < m; ++id) if (!chosen[id]) {
        auto [u, v, weight] = edges[id];
        if (u == v) continue;
        int removed = pathMaximum(u, v);
        ll candidateWeight = result.mstWeight - get<2>(edges[removed]) + weight;
        tuple<ll, int, int> candidate = {candidateWeight, id, removed};
        if (!result.alternative || candidate < *result.alternative)
            result.alternative = candidate;
    }
    return optional<Result>{move(result)};
}
\end{lstlisting}
% Для неориентированного взвешенного мультиграфа возвращается
% \texttt{optional} локального результата: \texttt{mstWeight},
% \texttt{mstEdges} (ID выбранного минимального остова), \texttt{alternative}.
% Последнее поле --- \texttt{optional<tuple<ll,int,int>>}: стоимость лучшего
% другого остова, ID добавляемого и удаляемого ребра. Равная стоимость
% разрешена: отличается именно множество ID. При несвязности или $n=0$
% нет результата; при единственном остове нет альтернативы. Для $n=1$
% остов пустой, петли не дают альтернативы. Время
% $O(E\log(E+1)+(V+E)\log(V+1))$, память $O(E+V\log(V+1))$.
% Краскал выбирает первый остов, двоичные подъёмы находят максимальное
% ребро на пути. Для каждой нетривиальной хорды заменяется одно такое ребро.
% Все частичные суммы весов остова, разности \texttt{mstWeight-removedWeight}
% и последующие суммы с добавляемым весом должны помещаться в \texttt{ll}.
% Маркером отсутствующего ребра служит ID $-1$, поэтому допустим даже вес
% \texttt{LLONG\_MIN}. Размеры таблиц подъёмов помещаются в \texttt{int}.
%
% Для строго большего веса недостаточно отбрасывать равные кандидаты:
% нужно хранить два наибольших различных веса на пути. Здесь такая задача
% не решается намеренно. Возвращённая замена позволяет восстановить второй
% остов без хранения ещё одного полного списка рёбер.

\fi
